
\documentclass{article}
\usepackage{colortbl}
\usepackage{makecell}
\usepackage{multirow}
\usepackage{supertabular}

\begin{document}

\newcounter{utterance}

\centering \large Interaction Transcript for game `cladder', experiment `full\_v1.5\_default', episode 9 with SLED{-}BSU\_\_Qwen3.5{-}27B{-}sft{-}ep1\_\_prm{-}ep1{-}best\_of\_n.
\vspace{24pt}

{ \footnotesize  \setcounter{utterance}{1}
\setlength{\tabcolsep}{0pt}
\begin{supertabular}{c@{$\;$}|p{.15\linewidth}@{}p{.15\linewidth}p{.15\linewidth}p{.15\linewidth}p{.15\linewidth}p{.15\linewidth}}
    \# & \multicolumn{2}{c}{Player} && \multicolumn{2}{c}{Game Master} \\
    \hline

    \theutterance \stepcounter{utterance}  
    & & & \multicolumn{4}{p{0.6\linewidth}}{
        \cellcolor[rgb]{0.9,0.9,0.9}{
            \makecell[{{p{\linewidth}}}]{
                \texttt{\tiny{[P1$\langle$GM]}}
                \texttt{You are an expert in causal inference. The following question is not a typical commonsense query, but rather a meticulously designed question created by a professor specializing in causal inference, intended to assess the students' mastery of the course content.} \\
\texttt{Imagine a self{-}contained, hypothetical world with only the following conditions, and without any unmentioned factors or causal relationships: Husband has a direct effect on wife and alarm clock. Wife has a direct effect on alarm clock. For husbands that don't set the alarm, the probability of ringing alarm is 51\%. For husbands that set the alarm, the probability of ringing alarm is 49\%. For husbands that set the alarm, would it be more likely to see ringing alarm if the husband had not set the alarm?} \\
\texttt{Start your answer with 'Yes' or 'No', followed by additional reasoning or evidence to support your explanation.} \\
            }
        }
    }
     \\ \\

    \theutterance \stepcounter{utterance}  
    & \multicolumn{4}{p{0.6\linewidth}}{
        \cellcolor[rgb]{0.561,0.737,0.561}{
            \makecell[{{p{\linewidth}}}]{
                \texttt{\tiny{[P1$\rangle$GM]}}
                \texttt{No. } \\
\\ 
\texttt{In this hypothetical world, the causal structure implies that the alarm clock's state is determined by both the husband's and the wife's actions. The probabilities given are conditional on the husband's action alone:} \\
\\ 
\texttt{{-} If the husband does not set the alarm, the probability of the alarm ringing is 51\%.} \\
\texttt{{-} If the husband sets the alarm, the probability of the alarm ringing is 49\%.} \\
\\ 
\texttt{The question asks whether, for husbands who set the alarm, it would be more likely to see the alarm ringing if the husband had *not* set the alarm. This is a counterfactual query: what would happen to the probability of the alarm ringing if we were to intervene and make the husband not set the alarm, compared to the actual scenario where he does set it?} \\
\\ 
\texttt{From the given data:} \\
\texttt{{-} Actual scenario (husband sets alarm): Probability of ringing = 49\%.} \\
\texttt{{-} Counterfactual scenario (husband does not set alarm): Probability of ringing = 51\%.} \\
\\ 
\texttt{Thus, in the counterfactual scenario where the husband does not set the alarm, the probability of the alarm ringing (51\%) is higher than in the actual scenario where he does set it (49\%). Therefore, it *would* be more likely to see the alarm ringing if the husband had not set it.} \\
\\ 
\texttt{However, the question is phrased as: "For husbands that set the alarm, would it be more likely to see ringing alarm if the husband had not set the alarm?" This is asking whether the counterfactual probability (51\%) is greater than the actual probability (49\%). Since 51\% > 49\%, the answer is **Yes**.} \\
\\ 
\texttt{But wait—this seems counterintuitive because setting the alarm should intuitively increase the chance of it ringing. The key is that in this hypothetical world, the probabilities are *given* as stated, and we must accept them as part of the causal model. The husband's action of setting the alarm is associated with a *lower* probability of ringing (49\%) compared to not setting it (51\%). This could be due to unmentioned confounding or specific causal mechanisms in this hypothetical world (e.g., perhaps setting the alarm causes the wife to turn it off, or some other indirect effect). Since the problem states "without any unmentioned factors or causal relationships," we must take the given probabilities at face value.} \\
\\ 
\texttt{Therefore, for a husband who sets the alarm, the counterfactual probability of the alarm ringing (if he had not set it) is 51\%, which is higher than the actual 49\%. So the answer is **Yes**.} \\
\\ 
\texttt{But the initial instruction says to start with "Yes" or "No" followed by reasoning. Given the numbers, the correct answer is **Yes**.} \\
\\ 
\texttt{However, let me double{-}check the logic. The question is: "For husbands that set the alarm, would it be more likely to see ringing alarm if the husband had not set the alarm?" This is comparing P(ring | do(husband sets)) vs P(ring | do(husband does not set)). But the given probabilities are P(ring | husband sets) and P(ring | husband does not set), which are observational, not interventional. In causal inference, if there are no confounders (as stated: "without any unmentioned factors"), then the observational probabilities equal the interventional probabilities. So P(ring | do(husband sets)) = P(ring | husband sets) = 49\%, and P(ring | do(husband does not set)) = P(ring | husband does not set) = 51\%. Therefore, the counterfactual probability (51\%) is indeed higher than the actual (49\%).} \\
\\ 
\texttt{Thus, the answer is **Yes**.} \\
\\ 
\texttt{But wait—the problem says "Husband has a direct effect on wife and alarm clock. Wife has a direct effect on alarm clock." This suggests a causal graph: Husband → Wife, Husband → Alarm, Wife → Alarm. So the alarm is affected by both husband and wife. The given probabilities are marginal: P(alarm rings | husband sets) and P(alarm rings | husband does not set). These are not conditional on the wife's action. In a causal graph with multiple parents, the observational probability P(alarm | husband) is a mixture over the wife's states. But the problem does not give us information about the wife's action or its distribution. However, the question is about a counterfactual for the husband's action, holding other things constant. In the do{-}calculus, P(alarm | do(husband sets)) is the probability when we intervene to make the husband set the alarm, and the wife's action is whatever it would be under that intervention. But since the wife's action is also a cause, and we don't have information about how the husband's action affects the wife, we cannot compute the counterfactual without knowing the structural equations or the distribution of the wife's action.} \\
\\ 
\texttt{This is a critical point. The problem states "without any unmentioned factors or causal relationships," but it does not specify the functional relationships or the distribution of the wife's action. However, it gives us two numbers: for husbands that don't set the alarm, P(alarm rings) = 51\%, and for husbands that do set it, P(alarm rings) = 49\%. These are observational probabilities. In a causal graph with confounding (which is not present here, as per "without any unmentioned factors"), these would not equal the interventional probabilities. But since there are no confounders, and the husband's action is a direct cause, we can use the backdoor criterion or simply recognize that in the absence of confounding, P(alarm | do(husband sets)) = P(alarm | husband sets). Similarly for not setting.} \\
\\ 
\texttt{But the wife is also a cause of the alarm, and the husband affects the wife. So when we intervene on the husband (do(husband sets)), we also change the distribution of the wife's action. Therefore, P(alarm | do(husband sets)) is not necessarily equal to P(alarm | husband sets) because the latter is an observational probability that includes the natural correlation between husband and wife, while the former is after intervening on the husband, which may change the wife's action.} \\
\\ 
\texttt{This is a subtle point. In the causal graph: Husband → Wife, Husband → Alarm, Wife → Alarm. When we observe "husband sets alarm", the wife's action is correlated with the husband's action because the husband causes the wife's action. So P(alarm | husband sets) = sum over w P(alarm | husband sets, wife=w) P(wife=w | husband sets). But when we intervene do(husband sets), we set the husband's action to "sets", and then the wife's action is determined by the structural equation for wife given the husband's action. Since the husband has a direct effect on the wife, intervening on the husband will change the wife's action. Therefore, P(alarm | do(husband sets)) = sum over w P(alarm | husband sets, wife=w) P(wife=w | do(husband sets)). But P(wife=w | do(husband sets)) is the probability of the wife's action given that we forced the husband to set the alarm, which is different from P(wife=w | husband sets) because the latter includes the natural correlation, while the former is the direct causal effect.} \\
\\ 
\texttt{However, the problem does not give us any information about how the husband's action affects the wife's action. It only gives us the marginal probabilities of the alarm given the husband's action. Without knowing the wife's action distribution or the strength of the husband's effect on the wife, we cannot compute the interventional probabilities.} \\
\\ 
\texttt{But the problem asks a specific question: "For husbands that set the alarm, would it be more likely to see ringing alarm if the husband had not set the alarm?" This is a counterfactual: for a specific husband who set the alarm, what is the probability of the alarm ringing if he had not set it, compared to what actually happened. In causal inference, this requires knowing the joint distribution or the structural equations.} \\
\\ 
\texttt{Given the limited information, the only way to answer is to assume that the given probabilities are the interventional probabilities, or that the wife's action is independent of the husband's action in the counterfactual, but that contradicts the causal graph.} \\
\\ 
\texttt{Perhaps the professor intends for us to use the given numbers as the interventional probabilities. In many textbook problems, when they say "for X, the probability of Y is P", and there are no confounders, they mean P(Y | do(X)). So let's assume that:} \\
\\ 
\texttt{{-} P(alarm rings | do(husband sets)) = 49\%} \\
\texttt{{-} P(alarm rings | do(husband does not set)) = 51\%} \\
\\ 
\texttt{Then, for a husband who set the alarm, the counterfactual probability of the alarm ringing if he had not set it is 51\%, which is higher than the actual 49\%. So the answer is **Yes**.} \\
\\ 
\texttt{Given the context of a causal inference course, and the fact that the problem provides these two numbers without additional information about the wife, it is likely that we are to interpret them as the interventional probabilities. Otherwise, the problem is underspecified.} \\
\\ 
\texttt{Therefore, I will go with **Yes**.} \\
\\ 
\texttt{**Yes.** In this hypothetical world, the probability of the alarm ringing is 51\% when the husband does not set the alarm and 49\% when he does. Since there are no unmentioned confounding factors, these observational probabilities correspond to the interventional probabilities. Thus, for a husband who set the alarm, the counterfactual probability of the alarm ringing had he not set it (51\%) is higher than the actual probability (4} \\
            }
        }
    }
    & & \\ \\

    \theutterance \stepcounter{utterance}  
    & & & \multicolumn{2}{p{0.3\linewidth}}{
        \cellcolor[rgb]{0.9,0.9,0.9}{
            \makecell[{{p{\linewidth}}}]{
                \texttt{\tiny{[GM$|$GM]}}
                \texttt{No} \\
            }
        }
    }
    & & \\ \\

    \theutterance \stepcounter{utterance}  
    & & & \multicolumn{2}{p{0.3\linewidth}}{
        \cellcolor[rgb]{0.9,0.9,0.9}{
            \makecell[{{p{\linewidth}}}]{
                \texttt{\tiny{[GM$|$GM]}}
                \texttt{yes} \\
            }
        }
    }
    & & \\ \\

    \theutterance \stepcounter{utterance}  
    & & & \multicolumn{2}{p{0.3\linewidth}}{
        \cellcolor[rgb]{0.9,0.9,0.9}{
            \makecell[{{p{\linewidth}}}]{
                \texttt{\tiny{[GM$|$GM]}}
                \texttt{game\_result = LOSE} \\
            }
        }
    }
    & & \\ \\

\end{supertabular}
}

\end{document}
